% ---- pg-285 （书页 273） ----

\begin{mvtable}
\phantomsection\label{tab:6.16}\mvtabcap{表 6.16\quad 骨骼中无机物含量(1 年后)}
\begin{tabular}{crrrrrr}
\toprule
被测对象 & 主桡骨 & 桡骨 & 主肱骨 & 肱骨 & 主尺骨 & 尺骨\\\midrule
1 & 1.027 & 1.051 & 2.268 & 2.246 & 0.869 & 0.964\\
2 & 0.857 & 0.817 & 1.718 & 1.710 & 0.602 & 0.689\\
3 & 0.875 & 0.880 & 1.953 & 1.756 & 0.765 & 0.738\\
4 & 0.873 & 0.698 & 1.668 & 1.443 & 0.761 & 0.698\\
5 & 0.811 & 0.813 & 1.643 & 1.661 & 0.551 & 0.619\\
6 & 0.640 & 0.734 & 1.396 & 1.378 & 0.753 & 0.515\\
7 & 0.947 & 0.865 & 1.851 & 1.686 & 0.708 & 0.787\\
8 & 0.886 & 0.806 & 1.742 & 1.815 & 0.687 & 0.715\\
9 & 0.991 & 0.923 & 1.931 & 1.776 & 0.844 & 0.656\\
10 & 0.977 & 0.925 & 1.933 & 2.106 & 0.869 & 0.789\\
11 & 0.825 & 0.826 & 1.609 & 1.651 & 0.654 & 0.726\\
12 & 0.851 & 0.765 & 2.352 & 1.980 & 0.692 & 0.526\\
13 & 0.770 & 0.730 & 1.470 & 1.420 & 0.670 & 0.580\\
14 & 0.912 & 0.875 & 1.846 & 1.809 & 0.823 & 0.773\\
15 & 0.905 & 0.826 & 1.842 & 1.579 & 0.746 & 0.729\\
16 & 0.756 & 0.727 & 1.747 & 1.860 & 0.656 & 0.506\\
17 & 0.765 & 0.764 & 1.923 & 1.941 & 0.693 & 0.740\\
18 & 0.932 & 0.914 & 2.190 & 1.997 & 0.883 & 0.785\\
19 & 0.843 & 0.782 & 1.242 & 1.228 & 0.577 & 0.627\\
20 & 0.879 & 0.906 & 2.164 & 1.999 & 0.802 & 0.769\\
21 & 0.673 & 0.537 & 1.573 & 1.330 & 0.540 & 0.498\\
22 & 0.949 & 0.900 & 2.130 & 2.159 & 0.804 & 0.779\\
23 & 0.463 & 0.637 & 1.041 & 1.265 & 0.570 & 0.634\\
24 & 0.776 & 0.743 & 1.442 & 1.411 & 0.585 & 0.640\\
\bottomrule
\end{tabular}
\end{mvtable}

\vspace{2pt}
\noindent\rule{\linewidth}{0.4pt}

\noindent\small 资料来源: 数据经 Everett Smith 许可使用.

\vspace{1em}
\noindent (a) 在 $\alpha=0.05$ 下检验上述假设.

\noindent (b) 构造均值差的 95\% 联合置信区间.

\noindent (c) 构造 95\% 的庞弗罗尼联合置信区间, 并与 (b) 中的置信区间进行比较.

\noindent 6.31\quad 花生是美国南部部分地区的重要农作物. 为改进花生的种植方法, 农作物学家们对若干个变量不断地比较花生品种, \hyperref[tab:6.17]{表 6.17} 中给出了一个双因子实验的结果数据. 在两块土地(1,2)上交替种植三个品种的花生(品种 5,6 和 8), 并且测量能代表产品的三个变量和两个重要的谷物等级指标. 三个变量分别是

\noindent $X_1=$ 产量

\noindent $X_2=$ 花生仁的成熟程度(250 克花生的果仁重量, 单位克)

\noindent $X_3=$ 种子大小(100 颗种子的重量, 单位为克)

\noindent 实验重复两次.

\noindent (a) 利用\hyperref[tab:6.17]{表 6.17} 中的数据, 完成双因子 MANOVA. 检验地块效应、品种效应和地块—品种交互效应的显著性(取 $\alpha=0.05$).

\noindent (b) 分析残差, 通常的 MANOVA 假定是否满足? 试作讨论.
